\begin{textsample}{POS Dim 3 – human – Score 5.00 – rj/rj\_inf\_e\_ef\_intro\_2\_p3.txt}  \label{ex:f3_pos_017}
Para atender a tais finalidades no âmbito da educação escolar, a Carta Constitucional, no Artigo 210, já reconhece a necessidade de que sejam “ fixados conteúdos mínimos para o ensino fundamental, de maneira a assegurar formação básica \textbf{comum} e respeito aos valores culturais e artísticos, \textbf{nacionais} e regionais ” ( BRASIL, 1988 ).
Com base nesses marcos constitucionais, a Lei de \textbf{Diretrizes} e Bases da Educação Nacional ( LDB ), no Inciso IV de seu Artigo 9o, afirma que \textbf{cabe} à União :

% matched lemmas: caber, comum, diretriz, nacional
\end{textsample}
